\chapter*{РЕАЛИЗАЦИЯ}
\addcontentsline{toc}{chapter}{РЕАЛИЗАЦИЯ}

На листинге 1 приведена реализация функции генерации раундовых ключей для AES-256.\\
\includelistingpretty
    {task-01.cpp}
    {}
    {Генерация раундовых ключей}

Функция key\_expansion генерирует 60 32-битных слов раундовых ключей из исходного 256-битного ключа. Используются операции RotWord (циклический сдвиг), SubWord (подстановка через S-блок) и XOR с константами раунда Rcon.
На листинге 2 приведена реализация функции шифрования одного 128-битного блока.\\
\includelistingpretty
    {task-02.cpp}
    {}
    {Шифрование блока}

Функция encrypt\_block реализует основной алгоритм шифрования AES-256. Входной 16-байтовый блок преобразуется в матрицу состояния 4×4, затем последовательно применяются раундовые преобразования.
На листинге 3 приведена реализация основных операций AES: SubBytes и MixColumns.\\
\includelistingpretty
    {task-03.cpp}
    {}
    {Операции SubBytes и MixColumns}

Операция SubBytes выполняет нелинейную подстановку каждого байта через S-блок. Операция MixColumns выполняет линейное преобразование столбцов матрицы состояния с использованием умножения в поле Галуа GF(2 ** 8).
На листинге 4 приведена реализация функции шифрования файла с блочной обработкой данных.\\
\includelistingpretty
    {task-04.cpp}
    {}
    {Шифрование файла}

Функция encrypt\_file выполняет шифрование произвольного файла. Файл читается блоками по 16 байт, каждый блок шифруется отдельно. В начало выходного файла записывается исходный размер для корректной расшифровки. Последний блок дополняется нулями до полного размера (zero padding).